% =============================================================================
% Chapter 5: Far-Field Force Program
% Input-ready fragment for the lean toy-model research proposal.
% This file intentionally contains no preamble or document environment.
% =============================================================================

\chapter{Far-Field Force Program}
\label{ch:far-field-force-program}

The far-field program asks what one common brane--bulk medium must do before the
project commits to a particular nonlinear particle core. At distances large
compared with a throat radius, and where appropriate large compared with the
brane thickness, a localized object can be represented provisionally by its
symmetry, multipole content, source parity, motion, and unresolved form factors.
The background response can then be used to determine the range, sign,
propagation, leakage, and unwanted channels that any eventual throat must
produce. This does not establish that a stable throat with the assumed source
exists. It establishes a capability contract that the later throat solution
must inherit.

The sectors are ordered so that each removes freedoms from the next. Gravity
first establishes a probe-independent source field, a fixed calibration, and
the role of distributed porous-brane return. Static electricity adds the reflection-odd
finite-slab response and the requirements for charge conservation and
universality. Magnetism asks whether motion of that same oriented source
produces the required transverse response without a second normalization.
Acceleration then tests whether static electricity, magnetism, and transverse
light are three limits of one current and one coupling structure. Passive
response, inertia, extra modes, and drag complete the far-field contract by
stating what a future dressed throat must do when it is placed in, moved
through, and accelerated within the same medium.

\section{Common method for the far-field sectors}

The far-field calculation begins on a declared throat-free slab background with
all fields needed by the selected conservative, relaxational, or mixed branch.
Let

\[
\bm{\Phi}
=
(\mathbf u_T,u_L,h_m,h_t,\delta n,\delta\theta,\delta\chi_B,\ldots)^T
\]

stand schematically for the complete constrained perturbation vector. The list
is not a commitment that every displayed coordinate remains independent after
constraints are solved. It is a reminder that the desired carrier cannot be
isolated from other variables merely by naming it. Static equilibrium response,
when a genuine equilibrium reduction exists, is governed by a Hessian and
static Green matrix,

\[
\mathcal H(\mathbf k),
\qquad
\mathcal G_{\rm stat}(\mathbf k)=\mathcal H^{-1}(\mathbf k),
\]

on the physical constrained subspace. Dynamic response is governed separately
by a retarded operator and its inverse,

\[
\mathcal O_R(\omega,\mathbf k),
\qquad
\mathcal G_R(\omega,\mathbf k)=\mathcal O_R^{-1}(\omega,\mathbf k).
\]

The static inverse determines long-range kernels only in an ensemble for which
an equilibrium interaction functional is legitimate. The retarded inverse
determines poles, phase lag, damping, outgoing radiation, and memory. A
conservative but intrinsically periodic source may require a cycle-averaged
Hamiltonian, Routhian, averaged action, or Noether-stress calculation rather
than a static free energy. A maintained dissipative source requires stress,
momentum-flux, conversion, and reservoir accounting and may not be assigned a
pair potential unless conservativity is independently established.

Before throats are solved, a localized source is represented by a vector
\(\mathbf{J}_a\) with its overall amplitude, internal profile, and
finite-thickness form factor left symbolic. The placeholder may carry known
transformation properties under reflection and motion, but it may not assume
that its magnitude is universal, that it is topologically conserved, or that
its static and moving normalizations agree. Those are later outputs. For an
oriented source it is useful to write

\[
\mathbf{J}_{\alpha,s}
=
 s\,\mathbf{J}^{\rm odd}_\alpha
+\mathbf{J}^{\rm even}_\alpha,
\qquad
s\in\{+1,-1\},
\]

where the odd and even vectors may each contain several coupled fields. This
parity decomposition does not determine range or force sign. It only states how
the source transforms on a reflection-symmetric background.

Each sector must perform both a desired-channel calculation and a complete
response audit. The desired calculation asks whether the common medium contains
the required low-wave-number stiffness, tensor structure, pole, source overlap,
and normalization. The audit asks what else the same source excites: additional
odd or even branches, the in-plane longitudinal mode, bulk sound, interface and
order modes, conversion and reservoir response, or the transverse light
continuum. A sector does not pass merely because its leading asymptotic term has
the right power law while an omitted branch is unstable, long ranged, strongly
radiative, or dissipative.

The expansion must also state its spatial regime. The strict leading far field
has

\[
R\gg \max(R_{\rm throat},H_{\rm br}),
\]

so the source is seen primarily through its lowest surviving projected
multipoles. An intermediate regime may satisfy \(R\gg R_{\rm throat}\) while
\(R\sim H_{\rm br}\), in which case finite-thickness form factors and heavier
normal branches can modify the kernel. The leading contract should identify the
asymptotic coefficient without pretending that it controls contact or
annihilation, while also recording the scale and operator data needed to compute
the crossover later. A result obtained only by replacing the two interfaces
with an infinitesimal sheet must be reopened if the neglected thickness is what
selects parity, screening, source normalization, or the sign of a correction.

Normalization must likewise be carried consistently from source to response.
For each desired carrier the calculation should specify the physical mode norm,
source overlap, observable convention, and flux or stress used to extract a
force or emitted power. This prevents the same eigenfunction from being
rescaled one way in the static sector and another way in the moving or
radiative sector. Dimensional agreement and the desired power law are necessary
checks, but they do not replace a common normalization tied to the parent
variables and boundary data.

The distinction between capability and realization remains strict. During this
chapter, a result may show that a source with a specified parity and continuity
law would generate an acceptable far field. After Candidate Parent System V1 is
frozen, the stable slab and complete spectrum must reproduce the same operators
and coefficients. After a throat is solved, its actual source vector, form
factor, throughput, conservation law, and motion must reproduce the provisional
inputs. A failure at either handoff reopens or rejects the corresponding claim.

Every sector therefore exports the same kinds of information to
\cref{ch:cross-sector-integration}: the fields and constraints it uses; the
static and retarded operator blocks it requires; coefficient signs and allowed
ranges; characteristic speeds and screening scales; source parity and
conservation requirements; normalization relations; prohibited extra channels;
and the regime in which the result is meant to hold. These outputs are
conjunctive. A correct range cannot compensate for a wrong sign, and a correct
static sign cannot compensate for unacceptable dynamics.

\section{Gravity}

\subsection{Behavior and proposed source-side material response}

The observer-facing target is an approximately central, attractive,
inverse-square gravity-like field over a declared local or intermediate regime.
The source should be independent of the species later used to probe it, and its
leading ordinary behavior should be unchanged when an otherwise identical
throat is reflected from the \(+w\) orientation to the \(-w\) orientation in
correspondingly reflected environmental data. At larger distances the field may
be modified by distributed porous-brane return, but screening, overshoot, or sign
reversal must be calculated explicitly rather than removed by convention.

The proposed material mechanism is broader than the slogan ``gravity is
inflow.'' A supported open throat supplies a localized, one-way constituent
drain. Its local region may generate pressure gradients, material acceleration,
ordered stress, ordered-to-bulk conversion, bulk compression, and interface
deformation. Material that crosses the throat mouth de-structures into the bulk
and does not return through that throat. Return is a separate, spatially
distributed process in which bulk material reorganizes into brane state across
the porous brane; in projected continuity it is represented by the positive
bulk-to-brane contribution to \(S_{\mathrm{leak}}\). Reservoir backreaction
couples the local drain and distributed return without making them one route or
one rate. The gravity calculation asks whether a definite combination of these
source-side and environmental quantities produces a three-dimensional field
before any test throat is introduced. Schematically,

\[
\vec{\mathfrak{g}}_{\rm src}
=
\mathcal G_{\rm src}
\!\left[
\mathbf J_n,\mathbf v_{\rm med},D_t\mathbf v_{\rm med},
n,P,\chi_B,\bm{\tau}_{\rm order},
\Gamma_B,\mathbf J_\chi,\mathcal R_{\rm porous};\mathcal B_\infty
\right].
\]

Here \(\mathcal R_{\rm porous}\) denotes the separate distributed porous-brane
\(S_{\mathrm{leak}}\) return profile and its coupled reservoir state. This
equation names candidate inputs; it does not assume that each survives the final
reduction or that the map is already known. In particular, the source is not
defined by how one selected probe accelerates. A reference convention may fix
the observer's acceleration units once,

\[
\boxed{
\mathbf g_{\rm eff}
=
\mathcal C_{\rm ref}\,
\vec{\mathfrak{g}}_{\rm src}
},
\]

but \(\mathcal C_{\rm ref}\) is then frozen. It may not depend on the later test
species or be adjusted with distance to hide a changing source. Raw material
velocity and the calibrated gravity observable are therefore not identified by
notation alone.

For a nearly spherical source, the desired local parameterization is

\[
\mathbf g_{\rm eff}(r)
=
-\frac{C_g\,\overline Q_g(r)}{r^2}\,\hat{\mathbf r},
\]

where \(\overline Q_g(r)\) is an orientation-even but radially signed enclosed
source coefficient after the source and environmental reflection test has
closed. The target covariance is

\[
\boxed{
\mathbf g_{\rm eff}
[\mathcal C_{\rm th}\mathcal T;
 \mathcal R_w\mathcal B_\infty]
=
\mathbf g_{\rm eff}
[\mathcal T;\mathcal B_\infty]
}.
\]

This is not the same as requiring equal coordinate-normal constituent flux.
Reflected outward-drain throats may have opposite global-\(w\) flux signs while
producing the same gravity source. Nor does orientation evenness establish
attraction. On the ordinary local branch the independent sign test is

\[
\overline Q_g(r_{\rm local})>0,
\]

in the displayed radial convention.

The separate, spatially distributed bulk-to-brane contribution to
\(S_{\mathrm{leak}}\) makes the gravity source potentially scale dependent. The
local one-way throat drain and its ordered-to-bulk conversion region may produce
a positive enclosed source, while return distributed across the porous brane
over a much larger region reduces it. A useful regime would
contain a matching window

\[
R_{\rm throat}\ll r\ll L_{\rm return},
\]

or another scale-independent near-to-far prescription, in which the local
coefficient is approximately constant. Beyond the return scale,
\(\overline Q_g(r)\) may approach zero, settle to another value, overshoot, or
change sign. The proposal does not select among those outcomes in advance. It
requires the global material and reservoir solution to determine them and to
show that any resulting departure is compatible with the intended application
of the analog.

\subsection{Required calculations and medium requirements}

The first calculation is a source-side reduction on a provisional localized
drain and conversion profile. It must begin from exact four-dimensional
constituent conservation and the selected order dynamics, project the complete
stress and transport response onto the brane, and identify the three-dimensional
quantity that has acceleration dimensions after one fixed calibration. Earlier
effective calculations coupling a localized sink to a separate distributed
porous-brane return motivate an inverse-square branch under stated assumptions,
but the current program must determine which combination of transport, pressure,
stress, ordered-to-bulk conversion, and the positive bulk-to-brane contribution
to \(S_{\mathrm{leak}}\) supplies the coefficient. It must also keep net throat
throughput, gross local order conversion, integrated distributed return, and the
final gravity-source coefficient distinct; none may be substituted for another
by definition.

The second calculation is the reflected-source comparison. Both orientations
of the provisional source must be embedded in reflected environmental data, and
the complete parent fields must be transformed. If the ordinary vacuum is
itself reflection symmetric, the two source fields should agree in that same
environment. If the upper and lower bulk-facing reservoirs differ, the
calculation should expose the resulting environmental splitting. A charge-
correlated gravity difference may be an unacceptable prediction, but it must
not be misidentified as an intrinsic failure of the throat reflection map
without first transforming the environment consistently.

The third calculation is the radial matching problem for the local one-way drain
and the separate distributed porous-brane \(S_{\mathrm{leak}}\) return. A local
source profile and a distributed return profile must be solved together far
enough to obtain \(\overline Q_g(r)\), not merely its near-source normalization. The
calculation must determine whether an inverse-square interval exists, how its
coefficient is extracted without an arbitrary radius choice, and where return
begins to screen or reshape it. The return profile must be universal,
environmentally calculable, or retained as an explicit departure. It cannot be
chosen separately for each throat species after the local source has been
computed.

The local-drain and distributed-return calculations must also close their
material and mechanical ledgers. A stationary source cannot remove ordered
material, momentum, or energy from the local slab while leaving the bulk and
porous-brane return system unchanged.
For a conservative branch, the relevant fluxes and stored energies must close
within the full medium. For a relaxational or mixed branch, the internal
reservoir, dissipation, and entropy production must be included. This does not
require the full cosmological solution during the far-field requirement pass,
but it does require enough accounting to rule out a gravity source maintained
by an implicit external pump.

The fourth calculation is dynamic. A time-dependent source must be coupled to
the full retarded background operator. Schematically,

\[
g_i(\omega,\mathbf k)
=
\mathcal R^g_{ij}(\omega,\mathbf k)
\mathcal S_j(\omega,\mathbf k).
\]

The poles of \(\mathcal R^g\) determine propagation, delay, phase lag, damping,
and radiation. The leading local compressional response may be associated with
a bulk-sound speed \(c_s\), but the complete gravity response may also involve
order, interface, return, and reservoir modes. The proposal therefore does not
assign gravity one propagation speed before the coupled operator is solved.
Every pole with nonzero source overlap must be recorded, including slow
relaxational tails and any vorticity or frame-dragging-like component of the
gravity flow.

These calculations impose several requirements on the common medium. It must
support a well-posed source-side reduction that is independent of later probe
species; a reflection-covariant ordinary source branch; a positive local radial
sign; an inverse-square matching regime or another controlled extraction of
local source strength; a calculable return profile; a stable dynamic response
with acceptable characteristic speeds and damping; and one frozen calibration.
The required source variables must also coexist with the light and electric
sectors. A constitutive choice that produces the desired gravity response while
destabilizing the slab, opening an unacceptable light continuum, or giving the
electric odd branch the wrong stiffness is not a gravity success for the common
program.

\subsection{Failure conditions and handoff}

The selected gravity branch fails if no probe-independent environmental field
can be constructed; if attraction appears only after inserting the response of
a chosen test object into the source definition; if reflected ordinary sources
do not yield the same gravity behavior under reflected symmetric data; if the
local radial sign is wrong; or if the inverse-square coefficient cannot be
extracted without an arbitrary observation radius. It also fails if distributed
return produces unacceptable species dependence, screening, overshoot, or sign
change; if the response contains ghosts, instabilities, or ill-posed modes within
the declared analog; if damping or delay is unacceptable; or if
one desired result requires a gravity-only coefficient not propagated through
the shared medium.

A narrower negative result need not reject the entire architecture. It may
reject one source functional, one order-dynamics branch, one return ansatz, or
one candidate calibration. The handoff to Chapter 6 must state the surviving
source variables, parity and sign requirements, matching and return scales,
retarded poles, characteristic speeds, calibration convention, and every
coefficient shared with the slab, electric, light, and reservoir sectors.

\section{Static electricity and charge}

\subsection{Oriented source, long-range carrier, and interaction sign}

The observer-facing static target is a signed long-range interaction whose
leading potential scales as \(1/R\) and whose leading force scales as
\(1/R^2\), with like orientations repelling and opposite orientations
attracting. The source sign is not assigned as a free scalar charge. It is
postulated to arise from the \(+w/-w\) orientation of a finite throat and is
represented in the far field as an oriented material boundary condition on the
two-interface slab.

For the graph-regime interfaces,

\[
h_m=\frac{h_++h_-}{2},
\qquad
h_t=\frac{h_+-h_-}{2},
\]

where \(h_m\) is reflection odd and \(h_t\) is reflection even. These are
leading collective coordinates, not guaranteed eigenmodes. The complete odd
and even response blocks may mix them with order, density, phase,
longitudinal, bulk, and reservoir variables. The static calculation must
therefore diagonalize the complete constrained Hessian rather than assume that
one bare height coordinate carries electricity.

A Coulomb-like equilibrium tail in three brane dimensions requires a healthy
odd eigenbranch with nonzero source overlap and leading \(k^2\) static
stiffness,

\[
\lambda_E(k)
=
\kappa_E k^2+O(k^4),
\qquad
\kappa_E>0.
\]

Its static Green function then has the target behavior

\[
G_E(0,k)
\sim
\frac{1}{\kappa_E k^2},
\qquad
G_E(R)
\sim
\frac{1}{4\pi\kappa_E R}.
\]

Reflection oddness is not enough. A gapped odd branch would be screened, while
a flexural branch with \(k^4\) static stiffness would have a different spatial
kernel even though it may also be gapless dynamically. Conversely, reflection
evenness does not guarantee short range. Every even branch with nonzero throat
coupling must be classified, and any intended thickness-like screened branch
must acquire its gap after all mixing and constraints are included.

The force sign must come from the physical source and boundary ensemble. A
healthy scalar field with an ordinary linear source term does not automatically
produce the required like-repels behavior. The model instead treats orientation
as a material boundary condition: same-oriented disturbances should compound
the departure from the preferred slab state, while opposite-oriented
disturbances should partially cancel it. To turn that picture into a result, an
equilibrium branch must compute the correct constrained or Legendre-transformed
on-shell energy with one-body self terms removed. A conservative periodic pair
must use a derived cycle-averaged functional or averaged stress. A driven branch
must obtain the force from the complete stress, momentum-flux, conversion, and
reservoir ledger.

It is useful to isolate the reflection-odd Coulomb monopole coefficient,

\[
C_{E,ab}(R)
=
g_a g_b+\delta C_{ab}^{\rm odd}(R),
\qquad
g_a,g_b\ge0,
\]

with all electric sign carried by \(s_a s_b\). The quantity
\(C_{E,ab}\) must not absorb even thickness forces, parity-mixed environmental
terms, or nonlinear core interactions. In an equilibrium description, the
leading odd contribution may be organized as

\[
E_{ab}^{\rm odd}(R)
=
 s_a s_b\,
\frac{g_a g_b+\delta C_{ab}^{\rm odd}(R)}{R}.
\]

At finite distance a positive coefficient is not by itself enough, because the
radial derivative of the correction also enters the force. The asymptotic
like-repels/opposite-attracts result therefore requires both the correction and
its scaled derivative to be subleading over the declared Coulomb regime. The
finite-thickness crossover and nonlinear core force remain later calculations,
developed in detail in the focused opposite-orientation companion rather than
inside this far-field chapter.

The static calculation must also inspect every additional long-range branch. A
second odd mode could alter charge universality or introduce an extra force. An
even gapless mode could survive after electric monopole cancellation and produce
a species-dependent neutral interaction. A parity-mixed term may appear if the
background or reservoirs break reflection symmetry. These are not automatically
fatal, but their range, sign, source overlap, and observational role must be
known before the electric sector can be handed to the common-medium synthesis.

\subsection{Universality, conservation, and dynamic status}

The long-range coefficient is not complete until its source dependence is
understood. If several independently solved unit-charge species later produce
an asymptotic coefficient matrix, removal of \(s_as_b\) should leave a leading
positive-semidefinite, approximately rank-one structure. That is the test that
allows one to assign a nonnegative one-throat magnitude \(g_a\). For ordinary
unit-charge branches, the further target is

\[
g_a\simeq g_0
\]

within the declared regime. Common slab thickness alone does not establish this.
The throat boundary amplitude, topology, support state, and normal source
profile must either be universal, quantized, topologically controlled, or
sufficiently decoupled from the leading electric eigenmode. Species-dependent
finite-size and polarization corrections may remain, but they must be
subleading and calculable.

The orientation label also requires a conservation structure. On a globally
two-sided, co-orientable brane, a candidate additive count is

\[
N_s=\sum_a s_a,
\]

or, more fundamentally, an oriented intersection number. A successful electric
sector must derive why continuous evolution conserves the relevant count, how
opposite orientations can be created or removed in pairs, and how the count
appears in a projected Gauss-like law. At the effective level the target is a
continuity equation

\[
\partial_t\rho_s+\nabla\cdot\mathbf J_s=0.
\]

This law is not established by writing it. It must descend from the permitted
parent evolution and remain valid through motion, deformation, pair creation or
annihilation, and the finite-thickness core description. If no global normal or
equivalent topological construction exists, the charge definition must be
reformulated rather than used unchanged.

The present architecture does not yet possess a demonstrated first-class
\(U(1)\) Gauss constraint. The intended result is therefore a material and
topological charge law whose projected far field is Maxwell-like in the stated
regime, not an assumption that exact electromagnetic gauge structure is already
present. The derivation must show which constraint or conservation identity
replaces the missing structure and what additional longitudinal, normal, or
material degrees of freedom remain observable. An additive count that is
conserved only in the static ansatz but fails during motion or topology change
would not close the charge sector.

The electric carrier's dynamical status is a separate gate. The static Hessian
may supply a \(k^2\) inverse without determining whether the corresponding
degree of freedom is propagating, constrained, relaxational, diffusive, or part
of a mixed response. The complete retarded odd block
\(\mathcal O_{R,{\rm odd}}(\omega,\mathbf k)\) must identify its poles, damping,
phase lag, continuum placement, and energy or passivity properties. With
ordinary local inertia, the Coulomb branch may have
\(\omega_E\simeq c_Ek\), but this is not assumed before diagonalization. A
constrained branch must show how a time-dependent conserved source is satisfied
without an omitted propagating degree of freedom. A relaxational branch must
close its internal reservoir and show that low-frequency damping, noise or
memory, and entropy production are acceptable.

This dynamic classification is important even before the radiation chapter.
A propagating electric mode may create an additional radiation channel. A slow
mode may produce Cherenkov-like emission or drag for ordinary motion. A
diffusive mode may generate persistent lag or attenuation. A mode placed near a
bulk or longitudinal continuum may leak. Conversely, a constrained response
may reduce radiation but impose a nontrivial instantaneous or elliptic closure
that must remain mathematically and materially consistent with the parent
system. Static success cannot postpone these questions indefinitely.

The electric handoff therefore includes more than \(\kappa_E\). It must record
the complete odd and even eigenvectors, source form factors, screening lengths,
continuum thresholds, leading factorization structure, allowed corrections,
conservation law, dynamic branch, characteristic speed or relaxation scale,
and all shared coefficients. The same \(g_a\), source profile, and current must
then be inherited by magnetism and radiation.

\subsection{Failure conditions and handoff}

The selected electric branch fails if the complete constrained system contains
no healthy odd eigenmode with nonzero source overlap and leading \(k^2\) static
stiffness; if the desired force sign appears only after choosing an unjustified
energy ensemble or inserting a sign by hand; if an additional odd or even
long-range mode produces unacceptable forces; or if the full static operator is
unbounded, ghostlike, or unstable. It also fails if the leading coupling does
not factorize sufficiently to define individual charge magnitudes, if ordinary
unit-charge species are unacceptably nonuniversal, or if the asymptotic
corrections and their derivatives are not genuinely subleading.

The charge program further fails if no conserved oriented count or current can
be defined, if charge depends on an inconsistent global orientation, if the
dynamic carrier has unacceptable poles, damping, leakage, or radiation, or if
the required result can be obtained only by assigning electricity a private
field or coefficient that is absent from light, magnetism, the throat, and the
common parent action. Failure of one candidate odd branch may leave another
permitted realization open; failure of every admitted route to a conserved,
healthy Coulomb carrier would reject the electric architecture of the model.

\section{Magnetism}

\subsection{Moving oriented source and transverse response}

Magnetism is the first direct test that static charge belongs to a dynamical
source rather than to an isolated boundary-value analogy. At low speed, motion
of the same oriented object should produce a sign-dependent transverse source
of schematic form

\[
\mathbf J_s\simeq \rho_s\mathbf V,
\qquad
\mathbf J_T\propto s\mathbf V.
\]

The target brane response is rotationally covariant and transverse, with the
leading tensor and distance dependence appropriate to a slowly moving localized
charge. In the simplest isotropic far field this means a magnetic-like response
whose direction is organized by the cross product of source velocity and
separation and whose leading magnitude has the corresponding inverse-square
field falloff. The exact coefficient and any medium-dependent corrections must
be derived from the shared retarded transverse operator rather than copied from
Maxwell theory.

The calculation should begin from the conserved oriented source obtained in the
static-electric package, expand a translated source profile to first order in
\(\mathbf V\), and project it onto every transverse background eigenmode. The
physical magnetic observable \(\mathbf B_{\rm eff}\) must be defined through a
fixed source and response convention, including its transformation under
orientation reversal, velocity reversal, spatial rotations, and time reversal.
The far-zone tensor, power law, phase, and finite-thickness form factor can then
be extracted from the retarded Green matrix. Existing exploratory work supports
the expected moving-source tensor structure, falloff, and velocity order under
stated assumptions, but the source normalization and sign remain conditional on
the unresolved electric boundary data and throat interior.

A moving source also perturbs the background in sign-independent ways. The
complete calculation must separate at least three effects:

\begin{enumerate}[leftmargin=2.2em,itemsep=0.35em]
  \item electromagnetic magnetism, which reverses with \(s\) and with
        \(\mathbf V\);
  \item inertial rebuilding, compression, and wake structure of the dressed
        medium, which is not principally electric-sign dependent; and
  \item velocity-dependent gravity, including any change in localized throat
        throughput or conversion, in the separate distributed
        \(S_{\mathrm{leak}}\) return profile, or in vorticity of the
        gravity-associated flow.
\end{enumerate}

The last effect may include a gravitomagnetic or frame-dragging-like component,
but it is not electromagnetic magnetism merely because both involve transverse
or vortical response. Their source parity, calibration, and coupling to probes
must remain distinct.

\subsection{Shared normalization and unwanted responses}

The decisive magnetic requirement is inheritance. If the electric far field
has a leading source magnitude \(g_a\), then the same \(g_a\), the same
orientation sign \(s_a\), the same normal source form factor, and the same
continuity law must determine the leading moving current. The magnetic
coefficient may depend on already-frozen transverse-mode normalization and
characteristic speed, but it may not be independently fitted. Schematically,

\[
\text{electric source data}
+\text{source velocity}
+\text{shared transverse response}
\longrightarrow
\mathbf B_{\rm eff}.
\]

This relation is stronger than obtaining the correct tensor form. Many media can
produce a velocity-dependent transverse field. The unification claim requires
the electric and magnetic normalizations, relative sign, and source dependence
to arise from one coupling structure. A source whose static response couples to
one eigenmode while its moving response requires an unrelated coefficient or a
different current has not closed the electromagnetic sector.

The unwanted-response audit must include excitation of the odd electric mode,
even thickness and order modes, the longitudinal branch, bulk sound, transverse
light, and conversion or reservoir modes. Uniform slow motion should not
continuously radiate below the applicable thresholds. Any stationary comoving
wake must be distinguished from an outgoing wave, and any dissipative phase lag
must be included in the force and power balance. The calculation must also test
whether the magnetic observable acquires unacceptable dependence on the
particle species, internal support state, direction relative to a preferred
medium frame, slab asymmetry, or environmental reservoir.

Because the target is Maxwell-like rather than exact Maxwell theory, the
calculation must also record any time-reversal or constitutive departure carried
by the magnetic observable. Such a departure is not automatically fatal in a
toy analog, but it must remain compatible with the static source law, the
radiation response, and the declared observer-facing regime. Calling a
transverse material swirl \(\mathbf B_{\rm eff}\) does not by itself establish
its transformation properties or its force on a later moving probe.

The magnetic branch fails if it has the wrong orientation or velocity parity,
wrong leading tensor or falloff, instability, or unacceptable extra channels.
It also fails as an electromagnetic-unification result if its normalization can
be obtained only by a new independent coupling. Its handoff to Chapter 6 is the
complete first-order moving-source kernel, tied normalization, transformation
properties, characteristic speed and phase behavior, and a separated inventory
of magnetic, inertial, and gravity-associated responses.

\section{Accelerating charge and electromagnetic radiation}

\subsection{One current across electricity, magnetism, and light}

Acceleration tests whether the static, moving, and radiative limits belong to
one source hierarchy,

\[
O(v^0):\ \text{static electricity},
\qquad
O(v):\ \text{magnetism},
\qquad
O(a):\ \text{transverse radiation}.
\]

The target is not simply that an accelerating defect makes some waves. The same
conserved oriented density and current, with the same leading one-throat
coupling magnitude and source profile, must generate all three orders. The
normalized transverse light modes obtained in
\cref{ch:light-research-program} provide the receiving states. The accelerating
source calculation must show how energy and momentum leave the dressed source
and enter both physical transverse polarizations with no separately chosen
radiation coefficient.

At the effective level the source obeys

\[
\partial_t\rho_s+\nabla\cdot\mathbf J_s=0,
\]

while its acceleration-dependent perturbation is projected through the complete
retarded response. For each normalized outgoing light mode \(\lambda=1,2\),
the far-zone amplitude is determined by the overlap of the source with that
mode's eigenvector, normal profile, and flux normalization. The exact formula is
candidate dependent, but the logic is fixed:

\[
(\rho_s,\mathbf J_s;g_a)
\xrightarrow{\ \mathcal G_R\ }
\{\mathcal A_1,\mathcal A_2,\mathcal A_{\rm extra}\}.
\]

The emitted transverse power, polarization pattern, angular dependence, and
frequency scaling must follow from those amplitudes and the same energy-momentum
normalization used for free light. A calculation that rescales the outgoing
modes or source separately from the light chapter would not test a shared
coupling.

The source-side energy balance must include reversible storage as well as
radiation. Acceleration can reconfigure the standing support mode, throat
geometry, near-field electric and gravity response, pressure distribution, and
conversion region. The external power should therefore be organized
schematically as

\[
P_{\rm ext}
=
\frac{dE_{\rm dressed}}{dt}
+P_{u_T}
+P_{\rm extra}
+P_{\rm diss},
\]

with momentum carried by the corresponding fields and reservoir. Reversible
near-field rebuilding contributes to the time derivative and later to inertia;
it must not be counted as radiation. Conversely, outgoing flux may not be
hidden inside a frequency-dependent inertial coefficient without being included
in the power ledger.

This stage also fixes the dynamic interpretation of the electric carrier. If it
is a propagating mode, its acceleration-induced radiation must be calculated.
If it is constrained, the constraint must remain compatible with the conserved
source and the retarded transverse response. If it is relaxational or mixed,
its phase lag, dissipation, internal reservoir, and possible stochastic partner
must be explicit. Static Coulomb behavior, magnetic response, and transverse
light emission cannot be declared unified while the main electric degree of
freedom remains dynamically unspecified.

\subsection{Additional radiative channels and energy--momentum transfer}

The accelerating source couples to the complete medium, not only to the desired
light branch. The calculation must separately determine

\[
P_{u_T},
\quad
P_{\rm odd},
\quad
P_{\rm even},
\quad
P_{u_L},
\quad
P_{\rm bulk},
\quad
P_\chi,
\]

where \(P_{u_T}\) includes both transverse light polarizations; the odd and even
terms cover the electric-normal and thickness/order sectors; \(P_{u_L}\) is
in-plane longitudinal emission; \(P_{\rm bulk}\) covers density or sound modes
in the surrounding phase; and \(P_\chi\) covers conversion or internal-reservoir
response. Additional branches discovered by the complete spectrum must be added
rather than folded into one residual loss number.

For every channel, the project must establish the receiving dispersion, source
Fourier support, mode-projector overlap, interface overlap, outgoing boundary
condition, energy or norm sign, and emitted flux. Kinematic intersection with a
continuum is not enough to prove radiation, but a desired symmetry suppression
must be derived rather than asserted. Harmonics and sidebands generated by an
intrinsically periodic supported throat may radiate even when the mean source is
slow. A dissipative branch may absorb power without a propagating far-zone wave.
Both effects belong to the closure test.

The desired ordinary regime requires transverse emission to dominate in the
specified electromagnetic process while every extra channel is absent,
symmetry suppressed, screened, threshold closed, or quantitatively acceptable.
The proposal does not require exact Maxwell theory. It permits calculable
longitudinal, bulk, order, or preferred-medium departures, but they must be
small enough in the declared regime and remain consequences of the one frozen
medium. An extra channel cannot be suppressed by reducing a coupling that the
static-electric or magnetic calculation already fixed unless the change is
propagated through all three sectors.

This stage fails as electromagnetic closure if the static, magnetic, and
radiative normalizations cannot be tied to one source; if acceleration does not
couple correctly to both transverse light polarizations; if energy and momentum
cannot be balanced; if the electric carrier's dynamic branch is inconsistent;
or if unavoidable extra radiation is unacceptable. A narrower failure may
reject one source profile, one speed ordering, or one candidate parent response.
The handoff must record the tied current and coupling, normalized transverse
emission kernel, complete channel powers, frequency and angular dependence,
retarded poles, and every speed or threshold that constrains the common medium.

\section{Passive response, inertia, extra modes, and drag}

\subsection{Active, passive, and inertial roles}

The far-field program must not use the word \emph{mass} as though one scalar
coefficient had already been derived. Active gravitational strength, passive
response, and inertial resistance are different measurements of the complete
object. They may ultimately obey universal low-energy ratios, but their equality
cannot be imposed by one normalization.

Active strength is a source-side quantity. Once the local signed gravity
coefficient, matching prescription, and calibration are fixed, an observer may
parameterize the local source by an active mass. That definition records the
strength of the derived environmental field; it does not identify the source
with trapped-wave energy, constituent rest mass, aperture throughput alone, or
any other single microscopic quantity. Return may make the effective active
source scale dependent outside the local matching regime.

Passive response belongs to a later test throat placed in the already calibrated
external field. In the low-frequency limit, write

\[
F_i^{\rm test}(\omega\!\to\!0)
=
\mathcal M_{ij}^{\rm passive}(0)
 g_{{\rm eff},j}+\cdots.
\]

The source field is held fixed while the response of different throat species
is calculated. On the ordinary stable branch, the symmetric quasistatic
response must be positive definite. A reference throat may set units, but the
calibration may not be altered for each later species to force equal
acceleration. Tensor, frequency, field-strength, and environmental corrections
must remain visible outside the declared universal regime.

Inertial response is the resistance encountered when the complete dressed
throat is accelerated. The relevant object includes the support mode, throat
shape, brane deformation, near-field flow, pressure and density response,
conversion region, and reservoir dressing. Two broad reversible contributions
are anticipated. The trapped standing pattern may need to retime its
counterpropagating components, shift phases and nodes, or alter its resonant
geometry. The surrounding material configuration must also be rebuilt as the
throat velocity changes. Energy and momentum stored in those adjustments can
produce an added-inertia-like response without identifying observed mass with a
sum of microscopic constituent masses.

For a closed reversible branch, the low-speed inertial tensor may be extracted
from the total momentum of the complete solution,

\[
M_{ij}^{\rm eff}
=
\left.
\frac{\partial P_i}{\partial v_j}
\right|_{\mathbf v=0},
\]

with all field and near-zone contributions included. For a relaxational or
mixed branch, the response is frequency dependent,

\[
F_i(\omega)
=
\left[
-\omega^2M_{ij}(\omega)
-i\omega\Gamma_{ij}(\omega)
+\cdots
\right]
\delta x_j(\omega).
\]

The real, reactive part describes stored momentum and effective inertia; the
absorptive part describes damping, conversion lag, heat, or reservoir
excitation. Both must follow from one closed retarded system. An unexcited
internal reservoir must be passive: a small periodic perturbation may deposit
zero or positive net work into the complete throat-plus-reservoir system, but it
may not extract unlimited work from an unprepared bath.

Only after positive ordinary active, passive, and inertial behavior is
established are universality ratios meaningful. The target low-energy tests have
the schematic form

\[
\frac{M_{\rm passive}}{M_{\rm inertial}}
\longrightarrow \eta_{\rm PI},
\qquad
\frac{M_{\rm active}^{\rm local}}{M_{\rm inertial}}
\longrightarrow \eta_{\rm AI}
\]

in a weak-field, low-velocity, low-frequency, homogeneous environment and a
valid local gravity matching regime. The constants may be normalized only after
the source and response conventions are frozen. Universality of a negative,
singular, unstable, or active response does not count as success. Nor would
these ratios alone establish the full geometric equivalence principle of GR;
they are the model's operational low-energy universality requirements.

The far-field stage can determine what source and response kernels are required,
which background modes contribute, and how species-dependent source residues
would enter. It cannot complete the mass test until moving supported-throat
families exist. Its handoff to the throat program is therefore explicit: each
solved species must later supply active source data, passive response, inertial
response, and their dependence on internal state without changing the common
background calibration.

\subsection{Complete mode inventory, thresholds, leakage, and friction}

A preferred material frame makes the complete spectrum part of the force
phenomenology. The model contains, or may contain after full diagonalization,
transverse guided light, in-plane longitudinal motion, odd electric/normal
response, even thickness response, bulk sound or density branches, order modes,
interface modes, and reservoir relaxation. Gravity may involve more than one of
these branches. Their characteristic speeds are not automatically equal:
\(c_\gamma\), the bulk sound speed \(c_s\), a possible electric speed \(c_E\),
and the longitudinal speed \(c_L\) arise from different combinations of common
coefficients. Cone alignment may emerge, remain approximate, or fail. Each
possibility changes leakage, radiation thresholds, and the observable
preferred-medium scope.

On a reversible homogeneous branch, uniform motion can emit into any coupled
mode satisfying the kinematic resonance condition. In an isotropic background,
a useful first diagnostic is

\[
v_c
=
\min_a\inf_{k>0}
\frac{\omega_a(k)}{k},
\]

where \(a\) runs over every physical branch with nonzero throat coupling. In an
anisotropic or defect background the threshold is directional and must use
\(\omega_a(\mathbf k)=\mathbf k\!\cdot\!\mathbf v\). A quadratic dispersion
\(\omega\propto k^2\) has vanishing \(\omega/k\) at long wavelength, so a
nonzero drag-free interval would then require symmetry suppression or a
vanishing source residue rather than merely a small velocity. A branch may also
be kinematically open yet dark because the relevant projection vanishes; the
source overlap and outgoing flux must be calculated.

Dissipative or mixed branches require an additional test. Absence of a
propagating wave does not imply absence of drag. The low-frequency friction,
order-relaxation, diffusion, and conversion kernels must vanish or remain within
the declared regime for uniform translation. A long-range static electric
response with a diffusive pole, for example, may leave a wake even when no
on-shell wave is emitted. The complete response should separate reversible
added inertia, acceleration radiation, and steady friction instead of fitting
them with one effective mass.

Motion may also change the source itself. For an isotropic throat, scalar
throughput and conversion corrections begin naturally at even powers of speed,
for example

\[
Q_A(v)
=
Q_{A0}
\left[
1+\alpha_A\frac{v^2}{c_*^2}+O(v^4)
\right],
\qquad
A\in\{n,\chi\}.
\]

The coefficients are outputs of the moving-throat family. Any resulting change
in the signed gravity source \(\overline Q_g(r;v)\) is a velocity-dependent
gravity effect, not inertial mass or electromagnetism. Likewise, velocity-
dependent electric coupling, polarization, slab thickness, or reservoir loading
would be a preferred-medium prediction that must be propagated through the
magnetic and radiation sectors.

The full audit therefore combines static range with dynamic risk. It asks which
modes are long ranged, which are screened, which are constrained, which can
carry energy away, which create low-frequency lag, and which couple to a
stationary, moving, or accelerating throat. It must include leakage from the
brane into bulk continua, conversion among transverse and longitudinal modes,
odd and even radiation, and reservoir dissipation. No unwanted branch may be
omitted because it is not part of the desired analogy.

The selected common-medium branch fails this stage if an ordinary stable throat
would possess a non-positive passive or inertial direction; if the retarded
response is active for an unexcited reservoir; if low-energy mass ratios are
unacceptably species dependent; if any unavoidable coupled mode gives a
vanishing or too-low critical velocity; if uniform motion has unacceptable
friction; or if the required characteristic-speed relations destroy light
confinement, Coulomb dynamics, or gravity response. A less severe result may
instead narrow the model to a preferred-medium regime with explicit bounded
departures. The scope of the conclusion depends on the calculated magnitude and
whether the conflict persists throughout the allowed common parameter region.

\section{Far-field requirement summary}

The force program does not end by assigning a separate pass label to gravity,
electricity, magnetism, radiation, and inertia. It ends by exporting one linked
set of source, response, and spectral requirements. The principal handoff is
summarized in \cref{tab:far-field-handoff}.

\begin{table}[htbp]
\centering
\caption{Principal far-field handoff to the common-medium synthesis.}
\label{tab:far-field-handoff}
\begin{tabularx}{\textwidth}{@{}L{0.20\textwidth}L{0.37\textwidth}Y@{}}
\toprule
\textbf{Category} & \textbf{Required output} & \textbf{Principal rejection mode} \\
\midrule
Gravity source
& Probe-independent source reduction; reflection-even signed coefficient;
  attractive inverse-square local regime; distributed porous-brane
  \(S_{\mathrm{leak}}\) return and matching map; one frozen calibration
& Wrong sign or parity, no usable local regime, species-dependent calibration,
  or unacceptable screening, overshoot, or retarded response. \\
Electric static sector
& Complete odd/even Hessian blocks; healthy source-coupled odd
  \(\kappa_Ek^2\) branch; correct physical force ensemble and sign; screened or
  bounded additional forces
& Wrong range or sign, ghost or instability, or an unavoidable extra long-range
  branch. \\
Charge universality
& Nonnegative one-throat magnitudes \(g_a\), bounded factorization corrections,
  conserved oriented count/current, and projected Gauss-like relation
& Independent normalization by species, no conserved current, or large
  nonfactorizing corrections. \\
Electric dynamics
& Propagating, relaxational, mixed, auxiliary, or constrained classification of
  the full odd retarded block, with poles, residues, damping, passivity, and
  zero-frequency closure
& Static Coulomb behavior necessarily entails unacceptable radiation, drag,
  dissipation, strong coupling, or constraint inconsistency. \\
Magnetism
& \(O(v)\) orientation-odd transverse source, required tensor and falloff, and
  normalization inherited from static charge and the light operator
& Independent magnetic coupling, wrong reversal law, or domination by an
  orientation-even wake. \\
Radiation
& \(O(a)\) transfer into both transverse light polarizations, shared current and
  coupling, radiation reaction, and channel-resolved energy--momentum balance
& Missing tied transverse emission, failed conservation, or excessive odd,
  even, longitudinal, bulk, or conversion radiation. \\
Mass response
& Positive active source, positive passive and inertial tensors, retarded
  passivity, and bounded low-energy species variation of the required ratios
& Wrong-sign or active response, singular inertia, or failed low-energy
  universality. \\
Speeds and drag
& Complete characteristic-speed, continuum, source-overlap, critical-velocity,
  leakage, and low-frequency friction map
& Unavoidable low-threshold emission, vacuum drag, cone incompatibility, or
  uncontrolled preferred-medium effects. \\
\bottomrule
\end{tabularx}
\end{table}

Several shared quantities appear repeatedly and must be represented only once in
Candidate Parent System V1: the slab thickness and background profiles; the
order branch and reservoir variables; the complete constrained static and
retarded operators; the transverse light normalization; odd and even source
profiles; the oriented current; the independent couplings; and the boundary and
environmental ensembles. The force requirements therefore constrain one another
before a throat is solved. A parameter that screens an unwanted even force may
change the photon helper. An electric kinetic term may open extra radiation. A
speed choice that reduces drag may move light or longitudinal modes into a bulk
continuum. A distributed porous-brane \(S_{\mathrm{leak}}\) return law that
produces the desired local gravity source may alter the reservoir state used by
every other sector.

The common-medium synthesis in \cref{ch:cross-sector-integration} will place
these outputs beside the light requirements and test their intersection. A
robust region must satisfy every hard structural gate and keep every bounded
departure inside its declared range at the same parameter point or throughout a
single declared region. A fragile point lying on a pole collision, continuum
threshold, sign change, or vanishing source residue is not equivalent to a
robust capability region, even if one numerical evaluation appears favorable.
An empty intersection rejects the candidate combination rather than authorizing
a private repair for the sector that failed.

Even a successful far-field intersection remains conditional until one frozen
parent system produces the stable slab and complete spectrum and then rederives
all of the responses above. It also remains conditional on the later throat
program: actual supported objects must realize the provisional source profiles,
conserved current, universal couplings, positive response, and acceptable motion.
The purpose of this chapter is to make that later solve decisive. The throat may
select among permitted source families, but it may not change what gravity,
electricity, magnetism, radiation, inertia, and drag were required to mean.
